\section{Precompile Architecture}


Zoo EVM activates 20 precompiles at genesis. These execute at native speed within the EVM, avoiding the 100--1000x overhead of implementing equivalent operations in Solidity bytecode.

\begin{table}[h]
\centering
\small
\begin{tabular}{llll}
\toprule
\textbf{Category} & \textbf{Precompile} & \textbf{Address} & \textbf{Standard} \\
\midrule
\multirow{5}{*}{Post-Quantum} & ML-DSA & \texttt{0x0200...06} & FIPS 204 \\
 & SLH-DSA & \texttt{0x0600...01} & FIPS 205 \\
 & ML-KEM & \texttt{0x0200...07} & FIPS 203 \\
 & PQCrypto & \texttt{0x...9003} & Composite \\
 & Blake3 & \texttt{0x0500...04} & BLAKE3 \\
\midrule
\multirow{3}{*}{Threshold Sig} & FROST & \texttt{0x0800...02} & RFC 9591 \\
 & CGGMP21 & \texttt{0x0800...03} & CGGMP21 \\
 & Corona & \texttt{0x0200...0B} & Ring-CT \\
\midrule
\multirow{3}{*}{Curves} & Ed25519 & \texttt{native} & RFC 8032 \\
 & secp256r1 & \texttt{native} & NIST P-256 \\
 & SR25519 & \texttt{native} & Ristretto \\
\midrule
\multirow{4}{*}{Privacy} & FHE & \texttt{0x0700...00} & CKKS/BGV \\
 & ZK & \texttt{0x0900...00} & Groth16 \\
 & HPKE & \texttt{0x...9200} & RFC 9180 \\
 & ECIES & \texttt{0x...9201} & IEEE 1363a \\
\midrule
Privacy & Ring & \texttt{0x...9202} & Borromean \\
\midrule
\multirow{2}{*}{DeFi} & DEX Pool & \texttt{0x...9010} & Uniswap v4 \\
 & Router & \texttt{0x...9012} & Swap Router \\
\midrule
\textbf{AI} & \textbf{AI Mining} & \texttt{0x0300...00} & \textbf{PoUC} \\
\midrule
Compute & Graph & \texttt{0x0500...10} & On-chain GNN \\
\bottomrule
\end{tabular}
\caption{Zoo EVM precompile suite. AI Mining (bold) is unique to Zoo---not present on the regulated EVM L1.}
\end{table}

\subsection{Post-Quantum Cryptography}

Zoo EVM implements all three NIST PQC standards finalized in 2024:

\textbf{ML-DSA} (FIPS 204): Module-Lattice Digital Signature Algorithm. Provides quantum-resistant digital signatures based on the hardness of Module Learning With Errors (MLWE). Zoo uses ML-DSA-65 (Level 3, 3.3 KB signatures) for validator attestations and governance votes.

\textbf{SLH-DSA} (FIPS 205): Stateless Hash-Based Digital Signature Algorithm. Provides signatures with minimal cryptographic assumptions (only hash function security). Larger signatures (17 KB for SLH-DSA-SHAKE-256f) but maximum conservative security. Used for long-term key certification.

\textbf{ML-KEM} (FIPS 203): Module-Lattice Key Encapsulation Mechanism. Provides quantum-resistant key exchange. Used in MPC keygen ceremonies and secure channel establishment between validators.

\subsection{AI Mining Precompile}

The AI Mining precompile at \texttt{0x0300...00} enables proof-of-useful-computation (PoUC):

\begin{enumerate}[nosep]
  \item \textbf{Task submission}: Research tasks (inference, training batches, evaluation) are posted on-chain with compute requirements and reward allocation.
  \item \textbf{Compute commitment}: Validators commit to executing tasks, staking ZOO as collateral.
  \item \textbf{Result verification}: Completed results are submitted with cryptographic proofs. For deterministic tasks (inference), bit-exact reproducibility is verified. For stochastic tasks (training), statistical validation against held-out test sets is used.
  \item \textbf{Reward distribution}: Verified compute earns ZOO rewards plus governance rights (ZIPs voting weight).
\end{enumerate}

This mechanism transforms idle validator compute into productive AI research capacity, aligning economic incentives with scientific output.

\subsection{Privacy Precompiles for ML}

Two precompiles enable privacy-preserving machine learning:

\textbf{FHE} (Fully Homomorphic Encryption): The CKKS/BGV precompile at \texttt{0x0700...00} enables computation on encrypted data. Applications include private model inference (user data never decrypted on-chain) and encrypted model parameter aggregation in federated learning.

\textbf{ZK} (Zero-Knowledge Proofs): The Groth16 precompile at \texttt{0x0900...00} enables succinct verification of computation. Applications include proving model accuracy without revealing model weights, verifying training data provenance, and privacy-preserving credential verification for DeSci access control.

